\documentclass[10pt,a4paper]{amsart}
\usepackage{amsmath,amssymb,amsthm,graphicx}
\usepackage{verbatim}
\newenvironment{reference}{\comment}{\endcomment}
%\newenvironment{reference}{}{}
\newenvironment{slogan}{\comment}{\endcomment}
\newenvironment{history}{\comment}{\endcomment}

% For commutative diagrams we use Xy-pic
\usepackage[all]{xy}

% We use 2cell for 2-commutative diagrams.
\xyoption{2cell}
% Theorem environments.
%
\theoremstyle{plain}
\newtheorem{theorem}[subsection]{Theorem}
\newtheorem{proposition}[subsection]{Proposition}
\newtheorem{lemma}[subsection]{Lemma}

\theoremstyle{definition}
\newtheorem{definition}[subsection]{Definition}
\newtheorem{example}[subsection]{Example}
\newtheorem{exercise}[subsection]{Exercise}
\newtheorem{situation}[subsection]{Situation}

\theoremstyle{remark}
\newtheorem{remark}[subsection]{Remark}
\newtheorem{remarks}[subsection]{Remarks}

\numberwithin{equation}{subsection}

% Macros
%
\def\lim{\mathop{\mathrm{lim}}\nolimits}
\def\colim{\mathop{\mathrm{colim}}\nolimits}
\def\Spec{\mathop{\mathrm{Spec}}}
\def\Hom{\mathop{\mathrm{Hom}}\nolimits}
\def\Ext{\mathop{\mathrm{Ext}}\nolimits}
\def\SheafHom{\mathop{\mathcal{H}\!\mathit{om}}\nolimits}
\def\SheafExt{\mathop{\mathcal{E}\!\mathit{xt}}\nolimits}
\def\Sch{\mathit{Sch}}
\def\Mor{\mathop{\mathrm{Mor}}\nolimits}
\def\Ob{\mathop{\mathrm{Ob}}\nolimits}
\def\Sh{\mathop{\mathit{Sh}}\nolimits}
\def\NL{\mathop{N\!L}\nolimits}
\def\CH{\mathop{\mathrm{CH}}\nolimits}
\def\proetale{{pro\text{-}\acute{e}tale}}
\def\etale{{\acute{e}tale}}
\def\QCoh{\mathit{QCoh}}
\def\Ker{\mathop{\mathrm{Ker}}}
\def\Im{\mathop{\mathrm{Im}}}
\def\Coker{\mathop{\mathrm{Coker}}}
\def\Coim{\mathop{\mathrm{Coim}}}

% Boxtimes
%
\DeclareMathSymbol{\boxtimes}{\mathbin}{AMSa}{"02}

%
% Macros for moduli stacks/spaces
%
\def\QCohstack{\mathcal{QC}\!\mathit{oh}}
\def\Cohstack{\mathcal{C}\!\mathit{oh}}
\def\Spacesstack{\mathcal{S}\!\mathit{paces}}
\def\Quotfunctor{\mathrm{Quot}}
\def\Hilbfunctor{\mathrm{Hilb}}
\def\Curvesstack{\mathcal{C}\!\mathit{urves}}
\def\Polarizedstack{\mathcal{P}\!\mathit{olarized}}
\def\Complexesstack{\mathcal{C}\!\mathit{omplexes}}
% \Pic is the operator that assigns to X its picard group, usage \Pic(X)
% \Picardstack_{X/B} denotes the Picard stack of X over B
% \Picardfunctor_{X/B} denotes the Picard functor of X over B
\def\Pic{\mathop{\mathrm{Pic}}\nolimits}
\def\Picardstack{\mathcal{P}\!\mathit{ic}}
\def\Picardfunctor{\mathrm{Pic}}
\def\Deformationcategory{\mathcal{D}\!\mathit{ef}}
\usepackage{cleveref}
\crefname{originallemma}{original lemma}{original lemmas}
\Crefname{originallemma}{Original lemma}{Original lemmas}
\makeatletter
\newlabel{more-morphisms-lemma-quasi-finite-quasi-section-meeting-nearby-open}{{05GT}{source}{}{sourcepointer-1}{}}
\newlabel{more-morphisms-lemma-quasi-finite-quasi-section-meeting-nearby-open@cref}{{[originallemma][0][]05GT}{[1][0][]source}{}{}{}}
\makeatother
\title{A local map of local domains essentially of finite type with finite residue extension that is not quasi-finite away from the closed point}
\author{The Stacks Project Authors}
\date{stacks-project:examples.tex:local-snapshot-20261012; GNU FDL 1.2 or later}
\begin{document}
\maketitle
\noindent\textit{Editorial note.} All delivered mathematical text is the author's own, copied from the pinned Stacks Project chapter file examples.tex; no statement or proof step was rewritten, reordered or removed.\par\medskip
\noindent\textit{Editorial note.} The counted proof is the author's complete argument: The counted proof is 'See the discussion above'; the spectrum analysis in the context (A has only (0) and m; C has only (0) over (0)) carries the argument; the final quasi-finite reading step is the author's 'desired example' conclusion, content verbatim.\par\medskip
\noindent\textit{Editorial note.} The label lemma-topology-finite-type names the counted statement, and the proof environment immediately following it in the source is delivered verbatim as its proof.\par\medskip
\noindent\textit{Editorial note.} The counted statement is carried by the author's own lemma environment in the chapter source.\par\medskip
\noindent\textit{Allowed standard prerequisites (editorial).}\begin{itemize}
\item Context span 3118-3159: A = \{a\_0 + a\_1 x + ... : a\_0 in k, a\_i in k((y))\} is local with maximal ideal m and no primes besides (0) and m; C similarly with k[y] constants, whose nonzero primes all contain p = \{a\_1 x + ... in C\} over m; B = C\_\{m\_C\} with m\_C = yC + p; the positive result for Noetherian A is More on Morphisms, lemma-quasi-finite-quasi-section-meeting-nearby-open
\item Localization of quasi-finite maps and the spectrum structure read off from the context
\end{itemize}\medskip
\ifdefined\hypertarget\hypertarget{sourcepointer-1}{}\fi
\noindent\textit{Source reference.} Original lemma 05GT: sources/stacks-chow/tags.txt, source lines 7072--7072. Prerequisite pointer invoked at 3181; it names a standard fact of a sibling chapter that the counted argument uses, and every step of the counted proof built on it is retained verbatim in the delivered body.\par\medskip

\noindent
Let $A \to B$ be a local map of local domains.
If $A$ is Noetherian, $A \to B$ is essentially of finite type,
and $A/\mathfrak m_A \subset B/\mathfrak m_B$ is finite then there
exists a prime $\mathfrak q \subset B$, $\mathfrak q \not = \mathfrak m_B$
such that $A \to B/\mathfrak q$ is the localization of a quasi-finite ring
map. See
More on Morphisms, Lemma
\ref{more-morphisms-lemma-quasi-finite-quasi-section-meeting-nearby-open}.

\medskip\noindent
In this section we give an example that shows this result is false $A$
is no longer Noetherian. Namely, let $k$ be a field and set
$$
A = \{a_0 + a_1 x + a_2 x^2 + \ldots \mid a_0 \in k, a_i \in k((y))
\text{ for }i\geq 1\}
$$
and
$$
C = \{a_0 + a_1 x + a_2 x^2 + \ldots \mid a_0 \in k[y], a_i \in k((y))
\text{ for }i\geq 1\}.
$$
The inclusion $A \to C$ is of finite type as $C$ is generated by $y$
over $A$. We claim that $A$ is a local ring with maximal ideal
$\mathfrak m = \{a_1 x + a_2 x^2 + \ldots \in A\}$ and no prime
ideals besides $(0)$ and $\mathfrak m$. Namely, an element
$f = a_0 + a_1 x + a_2 x^2 + \ldots$ of $A$ is invertible as soon as
$a_0 \not = 0$. If $\mathfrak q \subset A$ is a nonzero prime ideal,
and $f = a_i x^i + \ldots \in \mathfrak q$, then using properties of
power series one sees that for any $g \in k((y))$ the element
$g^{i + 1} x^{i + 1} \in \mathfrak q$, i.e., $gx \in \mathfrak q$.
This proves that $\mathfrak q = \mathfrak m$.

\medskip\noindent
As to the spectrum of the ring $C$, arguing in the same way as above
we see that any nonzero prime ideal contains the prime
$\mathfrak p = \{a_1 x + a_2 x^2 + \ldots \in C\}$ which lies
over $\mathfrak m$. Thus the only prime of $C$ which lies over
$(0)$ is $(0)$. Set $\mathfrak m_C = yC + \mathfrak p$ and
$B = C_{\mathfrak m_C}$. Then $A \to B$ is the desired example.


\begin{lemma}
\label{lemma-topology-finite-type}
There exists a local homomorphism $A \to B$ of local domains which is
essentially of finite type and such that $A/\mathfrak m_A \to B/\mathfrak m_B$
is finite such that for every prime
$\mathfrak q \not = \mathfrak m_B$ of $B$ the ring map
$A \to B/\mathfrak q$ is not the localization of a quasi-finite ring map.
\end{lemma}

\begin{proof}
See the discussion above.
\end{proof}

\end{document}
